\section{What remains open}\label{sec:open}

\paragraph{Which abstract recoveries can a tape realize cheaply?}
\Cref{thm:construction,thm:all-length} attain a rate loss
$O(\Delta^{-\alpha})$, including common partitions at every large
length with the all-subset tail bound. The unrestricted slow-gap
premise of \cref{cor:gap} is therefore false. Improving the seed or
the exponent may give stronger coding bounds, but that alone does
not decide (T).

The constructive question is whether a single uniform machine can
generate and evaluate useful recovery data, move it to the required
locations, and emit the transposed bits within linear time. The
abstract construction pays only for the number of labels. It does
not pay for accessing their coefficients or their values. Conversely,
a lower bound needs a restriction obeyed by the paid physical
computation that our unrestricted partitions need not satisfy.
Merely requiring a common partition, typical high-rate cells,
retention of the original coordinates, or all sufficiently large
lengths does not provide that distinction.

\paragraph{Which helpers can a physical computation use?}
\Cref{thm:native-tail} keeps the time at which each retained bit is
appended. This allows a prospective impossibility proof to choose
patches from a tape schedule and test them on its actual code. Fix a positive prediction advantage $\xi$ and a target
density $\alpha>0$. For each simulator batch size $s$, choose a per-patch
target bound $w_s$ before the dimension grows and use paired-distance
radius $2w_s-1$. Consider cells with the following property. Under the
uniform law on that cell, one fixed system scores at least an $\alpha$
fraction of its retained coordinates and predicts them from permitted
helpers with mean accuracy at least $1/2+\xi$. Its overlap is bounded by
a common $r_s=o(s)$. Must such cells carry a fraction of the uniform input
mass bounded away from zero as the dimension grows?

For a full successful cube at fixed $s$, the cells of deficit fraction
$O(1/s)$ carry all but at most $2^{-L/s}$ of the input mass. A positive
mass bound independent of the dimension would give a cell meeting both
conditions. On that cell \cref{cor:patch-prediction} would require overlap
of order $s$, contradicting the proposed upper bound at large $s$. If
targets or predictors are chosen separately for each input, the menu
test of \cref{cor:patch-menu} charges the number of systems used.
Establishing a prediction advantage on actual runs remains open.
Two independent fair helpers can determine their XOR exactly even though
either one alone gives chance accuracy.
Thus neither the abstract amplification nor the prediction inequality
supplies the missing physical-access theorem. A useful next result
would either derive the stated prediction advantage and overlap bound
from a linear-time schedule, or exhibit and account for a schedule that
avoids them while implementing the required recoveries.
